
\documentclass[12pt]{article}

% Layout.
\usepackage[top=1in, bottom=0.75in, left=1in, right=1in, headheight=1in, headsep=6pt]{geometry}

% Fonts.
\usepackage{mathptmx}
\usepackage[scaled=0.86]{helvet}
\renewcommand{\emph}[1]{\textsf{\textbf{#1}}}

% TiKZ.
\usepackage{tikz, pgfplots}
\usetikzlibrary{calc,trees,positioning,arrows,fit,shapes,calc}
\pgfplotsset{compat = newest}
\tikzset{
  jumpdot/.style={mark=*,solid},
  excl/.append style={jumpdot,fill=white},
  incl/.append style={jumpdot,fill=black},
}
 
\pgfplotsset{my style/.append style={axis x line=middle, axis y line=
middle, xlabel={$x$}, ylabel={$y$}, axis equal }}

% Misc packages.
\usepackage{amsmath,amssymb,latexsym,enumerate}
\usepackage{graphicx}
\usepackage{array}
\usepackage{xcolor}
\usepackage{multicol}

% Commands to set various header/footer components.
\makeatletter
\def\doctitle#1{\gdef\@doctitle{#1}}
\doctitle{Use {\tt\textbackslash doctitle\{MY LABEL\}}.}
\def\docdate#1{\gdef\@docdate{#1}}
\docdate{Use {\tt\textbackslash docdate\{MY DATE\}}.}
\def\doccourse#1{\gdef\@doccourse{#1}}
\let\@doccourse\@empty
\def\docscoring#1{\gdef\@docscoring{#1}}
\let\@docscoring\@empty
\def\docversion#1{\gdef\@docversion{#1}}
\let\@docversion\@empty
\makeatother

% Headers and footers layout.
\makeatletter
\usepackage{fancyhdr}
\pagestyle{fancy}
\fancyhf{} % Clears all headers/footers.
\lhead{\baselineskip 30pt
%\emph{\@doctitle\hfill\@docdate}
\emph{\@docdate\hfill\@doctitle}
\ifnum \value{page} > 1\relax\else\\
\emph{Name: \rule{3.5in}{1pt}\hfill \@docscoring}\fi}
\rfoot{\emph{\@docversion}}
\lfoot{\emph{\@doccourse}}
\cfoot{\emph{\thepage}}
\renewcommand{\headrulewidth}{0pt}%
\makeatother

% Paragraph spacing
\parindent 0pt
\parskip 6pt plus 1pt

% A problem is a section-like command. Use \problem{5} to
% start a problem worth 5 points.
\newcounter{probcount}
\newcounter{subprobcount}
\setcounter{probcount}{0}
\newcommand{\problem}[1]{%
\par
\addvspace{4pt}%
\setcounter{subprobcount}{0}%
\stepcounter{probcount}%
%\makebox[0pt][r]{\emph{\arabic{probcount}.}\hskip1ex}\emph{[#1 points]}\hskip1ex}
\makebox[0pt][r]{\emph{\arabic{probcount}.}\hskip1ex}\emph{[#1 \ifnum1=0#1\relax point\else points\fi]}\hskip1ex}
%\textbf{(#1 \ifnum1=0#1\relax thing\else things\fi
\newcommand{\thesubproblem}{\emph{\alph{subprobcount}.}}

% Subproblems are an enumerate-like environment with a consistent
% numbering scheme. 
% Use \begin{subproblems}\item...\item...\end{subproblems}
\newenvironment{subproblems}{%
\begin{enumerate}%
\setcounter{enumi}{\value{subprobcount}}%
\renewcommand{\theenumi}{\emph{\alph{enumi}}}}%
{\setcounter{subprobcount}{\value{enumi}}\end{enumerate}}

% Blanks for answers in normal and math mode.
\newcommand{\blank}[1]{\rule{#1}{0.75pt}}
\newcommand{\mblank}[1]{\underline{\hspace{#1}}}
\def\emptybox(#1,#2){\framebox{\parbox[c][#2]{#1}{\rule{0pt}{0pt}}}}

% Misc.
\renewcommand{\d}{\displaystyle}
\newcommand{\ds}{\displaystyle}
\def\bc{\begin{center}}
\def\ec{\end{center}}
\def\be{\begin{enumerate}}
\def\ee{\end{enumerate}}

\newcommand{\ans}{\rule{2 in}{.5pt}}
\newcommand{\bigans}[1]{\rule{#1 in}{.5pt}}

\doctitle{Math F251X: Quiz 6}
\docdate{Fall 2026}
\doccourse{UAF Calculus I}
\docversion{v-1}
\docscoring{\blank{0.8in} / 25}
\begin{document}
\textbf{Please circle your section:} \hfill 901 (Gossell)  \hfill 902 (Gossell) \hfill 903 (Graves-McCleary)\\

No aids (book, calculator, etc.) are permitted.  {\bf Show all work for full credit.}

%Problem
\problem{12} Find the derivatives of each of the following functions. Use correct derivative notation. You do not need to simplify your answers.
\begin{subproblems}
	\item $f(x)=\displaystyle e^x+x^e+e^e$
	\vfill
	\vfill
	\item $g(x)=\arctan(x^2)$
	\vfill
	\vfill
	\item $h(x)=\log_3(x)\arccos(e^x)$
	\vfill
	\vfill
	\vfill	
	\item $r(\theta)=\displaystyle\ln\left(\frac{\cos(\theta)}{\sqrt{2\theta+3}}\right)$ \hspace{1cm} (Hint:  You may want use log properties to simplify first.)
	\vfill
	\vfill
	\vfill
	\vfill
	\vfill
	\vfill
\end{subproblems}

\newpage

\problem{9} Part of the graph of $3^x-2\sin(y)=1$ is shown below.
\begin{multicols}{2}

\begin{subproblems}
	\item Sketch the \textbf{tangent line} to this graph at the point $(0,0)$.
	\vfill
	\item Use implicit differentiation to find $\displaystyle \frac{dy}{dx}$.
	\vfill
	\vfill
	\vfill
	\begin{center}
    \includegraphics[width=0.45\textwidth]{Q6Graph.png}
	\end{center}
\end{subproblems}
\end{multicols}
\vfill
\hspace{.3cm} \textbf{c.}  Use your answer in (b) to find the {\bf slope} of the tangent line to $(0,0)$.



\vfill

\problem{4} Considter the function $y=(3x+1)^x$.

\begin{subproblems}

\item Write down \textbf{your first step} in finding $\displaystyle \frac{dy}{dx}$.

\vfill

\item Choose the correct expression for $\displaystyle\frac{dy}{dx}$. ({\bf Circle one}.)

\begin{multicols}{2}
\begin{enumerate}[(I)]

\item $\displaystyle 3x(3x+1)^{x-1}$\\

\item $\displaystyle 3x(3x+1)^{x-1}\left(\frac{\ln(3x+1)}{3x+1}\right)$\\

\item $\displaystyle(3x+1)^x\left(\ln(3x+1)+\frac{x}{3x+1}\right)$\\

\item $\displaystyle (3x+1)^x\ln(3x+1)$\\

\item $\displaystyle(3x+1)^x\left(\ln(3x+1)+\frac{3x}{3x+1}\right)$\\

\item $\displaystyle(3x+1)^x + \frac{\ln(3x+1)}{3x+1}$\\

\end{enumerate}

\end{multicols}
\end{subproblems}




\end{document}
%%% Local Variables:
%%% mode: LaTeX
%%% TeX-master: t
%%% End:
